% Generated draft: requires mathematical and source review.
\documentclass[11pt]{article}
\usepackage[margin=25mm]{geometry}
\usepackage{fontspec}
\setmainfont{Latin Modern Roman}
\usepackage{amsmath,amssymb,mathtools,unicode-math}
\setmathfont{Latin Modern Math}
\usepackage{enumitem,needspace,xurl,hyperref}
\setlength{\parindent}{0pt}
\setlength{\parskip}{5pt}
\setlength{\emergencystretch}{3em}
\newcommand{\N}{\mathbb{N}}
\newcommand{\Z}{\mathbb{Z}}
\newcommand{\Q}{\mathbb{Q}}
\newcommand{\R}{\mathbb{R}}
\newcommand{\C}{\mathbb{C}}
\newcommand{\F}{\mathbb{F}}
\newcommand{\A}{\mathbb{A}}
\newcommand{\PP}{\mathbb{P}}
\newcommand{\E}{\mathbb{E}}
\newcommand{\eps}{\varepsilon}
\newcommand{\dd}{\,\mathrm{d}}
\begin{document}
\section*{UnsolvedMath ID: 30003438}
\textbf{Source problem code:} catalogue entry
\subsection{Definitions and mathematical statement}
Let
\[
S^2=\{x\in\R^3:\|x\|_2=1\}.
\]
The supplied record concerns the minimum pair-interaction energy of five particles on this sphere. Let \(\mathcal C_5\subseteq(S^2)^5\) denote the intended admissible configurations, let \(\mathcal P\) denote the intended class of pair potentials, and let \(\mathcal E_f(X)\) denote the intended total pair-interaction energy of a configuration \(X\in\mathcal C_5\) for \(f\in\mathcal P\). These are placeholders for definitions missing from the record. Write
\[
E_{\min}(f)=\inf_{X\in\mathcal C_5}\mathcal E_f(X).
\]
Denote the value of the source's second-level relaxation by \(E_2(f)\).

\textbf{Conjecture as supplied, with unresolved parameters.} Is this relaxation universally sharp, in the sense that
\[
\forall f\in\mathcal P,\qquad E_2(f)=E_{\min}(f)?
\]
The quantifier permits the minimizing configurations, if they exist, to depend on \(f\); it does not assert the existence of one configuration minimizing every potential.

\textbf{Unresolved qualifications.} The record does not define the relaxation \(E_2\), its feasible objects or constraints, or its objective normalization. It also does not specify the potential class \(\mathcal P\), the argument of a potential, or the formula and normalization of \(\mathcal E_f\). In particular, it does not say whether potentials act on chordal distance, squared chordal distance, or another quantity. The accompanying literature assessment mentions completely monotone potentials, but does not supply the interval, endpoint conditions, or the precise relation of that class to the statement's ``prescribed class.'' In standard terminology, complete monotonicity on an open interval \(I\) means that \(f\in C^\infty(I)\) and
\[
(-1)^k f^{(k)}(t)\geq0\qquad(k\in\Z_{\geq0},\ t\in I);
\]
this terminology alone does not determine the missing potential class here. Nor does the record specify whether coincident particles are excluded, how singular potentials are handled, or whether the energy infimum is attained. Finally, the assessment refers to a ``four-point bound,'' without defining it or establishing its identification with \(E_2\). Thus the displayed equality preserves the intended universal-sharpness question but is a schematic formulation requiring the missing source definitions before it becomes a fully specified mathematical conjecture.
\subsection{Short English statement}
For five particles on the unit sphere, the supplied conjecture asks whether the second-level relaxation gives exactly the least possible pair-interaction energy for every potential in the intended class. The relaxation, potential class, energy convention, and admissibility conditions are insufficiently specified in the supplied record. Its background mentions completely monotone potentials and a four-point bound, but their precise roles remain unresolved.
\subsection{Sources}
\begin{itemize}
\item Supplied local UnsolvedMath registry record, permanent ID \(30003438\), titled ``Universal Sharpness of Second-Level Five-Point Energy Relaxations,'' in the Oberwolfach Reports Open Problems collection. Exact source problem code: \texttt{catalogue entry}. The record gives proposed year 2017 and names no proposer.
\item Supplied report citation: \emph{Real Algebraic Geometry With a View Toward Moment Problems and Optimization} (2018), \emph{Oberwolfach Reports}, pp.~771--862, DOI 10.4171/owr/2017/14. Supplied link: \url{https://doi.org/10.4171/owr/2017/14}. The background describes Oberwolfach Reports as documenting mathematical workshops and their open-problem discussions. No problem-specific page or defining passage for \(E_2\) is supplied.
\item Supplied literature reference: de Laat, ``moment methods in energy minimization,'' \url{https://arxiv.org/abs/1503.07228}.
\item Supplied literature reference: Cohn--Kumar, ``universally optimal distributions,'' \url{https://arxiv.org/abs/math/0607446}.
\item Literature assessment dated exactly \texttt{2026-08-21}: the supplied curation labels the problem open and reports numerical semidefinite bounds and exact results for selected Riesz exponents for five points on the sphere, while describing universal sharpness of the four-point bound for every completely monotone potential as conjectural. This is supplied dated curation derived from the OWR record, not an independently verified assessment. The record itself requests independent verification of time-sensitive status and cited claims.
\item The registry labels the statement \texttt{corrected\_verified} and says mathematical content was preserved while notation, typography, and sentence structure were repaired for publication. These are supplied editorial assertions, not independent mathematical verification of this draft or its missing definitions.
\end{itemize}
\end{document}
